% R593: whole EGA I Chapter0 3.4 heading and3.4.1–3.
% OpenAI Codex, GPT-6.1 Sol, Ultra; no human review claimed.
\subsection{Imágenes directas de prehaces}
\label{subsection:0.3.4-fr}

\begin{env}[3.4.1]
\label{0.3.4.1-fr}
Sean $X$, $Y$ dos espacios topológicos y $\psi:X\to Y$ una aplicación
continua. Sea $\mathcal{F}$ un prehaz sobre $X$ con valores en una
categoría $\boldsymbol{K}$~; para todo abierto $U\subset Y$, sea
\[
  \mathcal{G}(U)=\mathcal{F}(\psi^{-1}(U)),
\]
y, si $U$, $V$ son dos partes abiertas de $Y$ tales que $U\subset V$, sea
$\rho_U^V$ el morfismo
\[
  \mathcal{F}(\psi^{-1}(V))\to\mathcal{F}(\psi^{-1}(U));
\]
es inmediato que los $\mathcal{G}(U)$ y los $\rho_U^V$ definen un
\emph{prehaz} sobre $Y$ con valores en $\boldsymbol{K}$, que se llama
la \emph{imagen directa de $\mathcal{F}$ por $\psi$} y se denota por
$\psi_*(\mathcal{F})$. Si $\mathcal{F}$ es un haz, se verifica inmediatamente
el axioma (F) para el prehaz $\mathcal{G}(U)$, por lo que
$\psi_*(\mathcal{F})$ es un haz.
\end{env}

\begin{env}[3.4.2]
\label{0.3.4.2-fr}
Sean $\mathcal{F}_1$, $\mathcal{F}_2$ dos prehaces sobre $X$ con valores
en $\boldsymbol{K}$, y sea $u:\mathcal{F}_1\to\mathcal{F}_2$ un morfismo.
Cuando $U$ recorre el conjunto de las partes abiertas de $Y$, la familia de
morfismos
\[
  u_{\psi^{-1}(U)}:
  \mathcal{F}_1(\psi^{-1}(U))\to\mathcal{F}_2(\psi^{-1}(U))
\]
satisface las condiciones de compatibilidad con los morfismos de restricción
y define, por tanto, un morfismo
\[
  \psi_*(u):\psi_*(\mathcal{F}_1)\to\psi_*(\mathcal{F}_2).
\]
Si $v:\mathcal{F}_2\to\mathcal{F}_3$ es un morfismo de $\mathcal{F}_2$ en
un tercer prehaz sobre $X$ con valores en $\boldsymbol{K}$, se tiene
\[
  \psi_*(v\circ u)=\psi_*(v)\circ\psi_*(u);
\]
en otras palabras, $\psi_*(\mathcal{F})$ es un \emph{funtor covariante} en
$\mathcal{F}$, de la categoría de los prehaces (resp. haces) sobre $X$ con
valores en $\boldsymbol{K}$ a la de los prehaces (resp. haces) sobre
$Y$ con valores en $\boldsymbol{K}$.
\end{env}

\begin{env}[3.4.3]
\label{0.3.4.3-fr}
Sean $Z$ un tercer espacio topológico y $\psi':Y\to Z$ una aplicación
continua, y sea $\psi''=\psi'\circ\psi$. Es claro que se tiene
\[
  \psi_*''(\mathcal{F})=\psi_*'(\psi_*(\mathcal{F}))
\]
para todo prehaz $\mathcal{F}$ sobre $X$ con valores en $\boldsymbol{K}$~; además,
para todo morfismo $u:\mathcal{F}\to\mathcal{G}$ entre tales prehaces,
se tiene
\[
  \psi_*''(u)=\psi_*'(\psi_*(u)).
\]
En otras palabras, $\psi_*''$ es el \emph{compuesto} de los funtores $\psi_*'$
y $\psi_*$, lo que puede escribirse
\[
  (\psi'\circ\psi)_*=\psi_*'\circ\psi_*.
\]

Además, para todo abierto $U$ de $Y$, la imagen por la restricción
$\psi|\psi^{-1}(U)$ del prehaz inducido
$\mathcal{F}|\psi^{-1}(U)$ no es otro que el prehaz inducido
$\psi_*(\mathcal{F})|U$.
\end{env}
